\subsection{Quality and recorded labor in offline policies}
\label{sec:labor-history}

Existing histories support lower target labor in the selection comparison, while acquiring histories changes the cost of adaptive stopping. At team size five, quality-and-time selection has mean $Q=93.926$ versus $93.708$ for historical-$Q$ selection across 126 contexts from 46 images in nine buildings. Their paired difference is $+0.218$ ($95\%$ building-bootstrap CI $[-0.538,0.869]$), and the geometric mean recorded-labor ratio is 0.813 ($[0.760,0.882]$).

Source-history stopping raises mean $Q$ from 85.189 to 86.054 on 21 images in five room groups. Its timing panel contains 11 images in three groups: target labor increases from 1,307.9 to 1,660.3 person-seconds, and including history acquisition gives a project-labor ratio of 2.115. The stopping rule improves observed quality and requires more work. Supplement D reports the full comparisons alongside Supplements C9, C13 and C14.
